\section{Fehlersuche anhand der Krumenstruktur}%
\label{sec:debugging-crumb-structure}

Die Krumenstruktur deines Brotes gibt Aufschluss darüber, wie gut
dein Fermentationsprozess verlaufen ist. Außerdem kannst du häufige Fehler
erkennen, die durch unsaubere Technik entstehen. Dieses Kapitel liefert dir Informationen,
mit denen du deinen Backprozess analysieren und verbessern kannst.

\begin{figure}
  \centering
  \includegraphics[width=\textwidth]{crumb-structures-book}
  \caption[Fehlersuche anhand der Krumenstruktur]{Eine schematische Darstellung
      verschiedener Krumenstrukturen und ihrer jeweiligen Ursachen. Die Krume
      des fertigen Brotes ist ein zentraler Aspekt, um mögliche Probleme bei
      der Fermentation oder der Backtechnik zu identifizieren.}%
  \label{fig:crumb-structures-book}
\end{figure}

\subsection{Perfekte Fermentation}

\begin{figure}
  \centering
  \includegraphics[width=\textwidth]{open-crumb}
  \caption[Perfekt fermentiertes Brot]{Das Brot hat eine eher offene Krume
      mit Bereichen, die eine Wabenstruktur aufweisen.}%
  \label{fig:open-crumb}
\end{figure}

Natürlich ist die perfekte Fermentation Ansichtssache und höchst subjektiv. Für
mich hat das perfekte Sauerteigbrot eine knusprige Kruste gepaart mit einer fluffigen,
etwas offenen Krume. Das ist die perfekte Balance unterschiedlicher Konsistenzen,
wenn man hineinbeißt.

Manche Menschen streben eine sehr offene Krume an, also große Lufteinschlüsse
(Alveolen). Es ist subjektiv, ob das der Brotstil ist, den du magst;
um ihn zu erreichen, musst du deinen Brotteig allerdings perfekt fermentieren.
Es braucht viel Können, sowohl in Bezug auf die Fermentation als auch auf die Technik,
um eine solche Krumenstruktur zu erzielen.

Persönlich mag ich~ein Brot dieser Art, nur mit einer etwas weniger wilden Krume.
Der Krumenstil, den ich~bevorzuge, heißt \emph{Wabenstruktur}. Er ist nicht zu offen, aber
gerade offen genug, um das Brot sehr fluffig zu machen. Um die zuvor erwähnte offene Krume zu erreichen,
muss man den Teig so wenig wie möglich berühren. Je mehr man mit dem
Teig arbeitet, desto mehr entgast man ihn. Übermäßiges Bearbeiten des Teigs
führt dazu, dass die Alveolen des Teigs zusammenfallen. Die Krume wird nicht so offen sein.
Deshalb funktioniert eine solche Krume am besten, wenn man nur einen
Laib auf einmal fermentiert. Normalerweise entgast man den Teig beim Rundwirken
automatisch ein wenig, wenn man ihn vorformen muss.
Wenn man diesen Schritt überspringt und den Teig direkt formt, erzielt man eine offenere Krume.
Als Faustregel gilt: Den Teig mindestens 1--2~Stunden vor dem Formen nicht mehr anfassen,
um eine möglichst offene Krume zu erzielen.

\begin{figure}
  \centering
  \includegraphics[width=\textwidth]{honeycomb}
  \caption[Wabenstruktur der Krume]{Ein Vollkorn-Sauerteigbrot mit einer fast
      ausschließlichen Wabenstruktur der Krume.}%
  \label{fig:honeycomb}
\end{figure}


Das wird allerdings problematisch, wenn man
mehrere Laibe gleichzeitig herstellen möchte. Das Vorformen ist unerlässlich, da man
den großen Hauptteig in kleinere Stücke teilen muss. Ohne das Vorformen
würde man viele ungleichmäßige Teigstücke erhalten. Diese Technik wird
auch bei Ciabattas angewandt. Sie werden typischerweise nicht geformt. Man schneidet den
Hauptteig nur in kleinere Stücke und versucht, den Teig so wenig wie möglich zu bearbeiten.
Durch das Vorformen nähern sich die Alveolen des Teigs eher einer Wabenstruktur an,
da große Lufteinschlüsse leicht zusammenfallen. Ähnlich wie bei der offenen Krumenstruktur
muss man auch hier den Fermentationsprozess perfekt treffen.
Eine zu lange Fermentation führt dazu, dass Gas beim Backen aus dem Teig entweicht.
Die Waben können das Gas nicht mehr halten. Fermentiert man zu kurz,
ist nicht genug Gas vorhanden, um die Strukturen aufzublähen. Für mich ist das der perfekte
Krumenstil. Als jemand, der Marmelade schätzt, fällt keine Marmelade durch eine Scheibe
dieses Brotes — im Gegensatz zur offenen Krume.

\subsection{Übergärt}%
\label{subsec:overfermented-dough}

\begin{figure}
  \centering
  \includegraphics[width=\textwidth]{fermented-too-long}
  \caption[Übergärtes Sauerteigbrot]{Ein relativ flacher Teig mit
      vielen winzigen Lufteinschlüssen.}%
  \label{fig:fermented-too-long}
\end{figure}

Wenn man den Teig zu lange fermentiert, beginnt das Enzym Protease,
das Gluten des Mehls abzubauen. Darüber hinaus verzehren die Bakterien das Gluten
in einem Prozess namens \emph{Proteolyse}~\cite{raffaella+di+cagno}.
Bäcker bezeichnen diesen Vorgang auch als \emph{Kleberabbau}.
Das Gluten, das normalerweise das \ch{CO2} aus
dem Fermentationsprozess der Mikroorganismen einschließt, kann das
Gas nicht mehr im Teig halten. Das Gas entweicht nach außen, was zu kleineren Alveolen in der Krume führt.
Das Brot selbst wird im Ofen sehr flach. Bäcker bezeichnen
dieses Brot oft als \emph{Fladen}. Der Ofentrieb ist vergleichbar
mit Brotteigen aus glutenarmem Mehl wie Einkorn.

Dein Brot wird viel Säure aufweisen, einen wirklich kräftigen, markanten Geschmack. Aus
geschmacklicher Sicht könnte es etwas zu sauer sein. Aus meinen eigenen Tests mit Familie und
Freunden (n=15--20) kann ich~sagen, dass diese Art von Brot typischerweise
weniger geschätzt wird. Allerdings mag ich~persönlich den herzhaften, kräftigen Geschmack sehr.
Er passt ausgezeichnet zu etwas
Süßem oder einer Suppe. Von der Konsistenz her ist das Brot nicht mehr so fluffig wie es sein könnte.
Die Krume könnte auch etwas gummiartig schmecken. Das liegt daran, dass sie von
den Bakterien stark abgebaut wurde. Darüber hinaus hat diese Art von Brot einen deutlich geringeren
Glutengehalt~\cite{raffaella+di+cagno} und ist nicht mehr mit rohem Mehl vergleichbar —
es ist ein vollständig fermentiertes Produkt.
Man kann es mit einem Blauschimmelkäse vergleichen, der fast laktosefrei ist.

Wenn man versucht, mit dem Teig zu arbeiten, fällt auf, dass er sich plötzlich
sehr klebrig anfühlt. Man kann den Teig nicht mehr richtig formen und bearbeiten. Wenn man
versucht, den Teig aus einem Gärkörbchen zu lösen, flacht der Teig stark ab. Außerdem
kann es in vielen Fällen passieren, dass der Teig am Gärkörbchen kleben bleibt. Als Anfänger beim Backen
habe ich~viel Reismehl in meinem Gärkörbchen verwendet, um die Oberfläche des Teigs stark zu trocknen.
So klebte der Teig nicht, obwohl er übergärt war. Wie sich jedoch
herausstellte, lag das Klebrigkeitsproblem an meinem mangelnden Verständnis des Fermentationsprozesses.
Heute verwende ich~nie Reismehl, außer für dekorative Einschnitte.
Eine korrekte Fermentation führt zu einem Teig, der nicht klebrig ist.

Wenn du während eines Dehnens und Faltens oder beim Formen merkst, dass dein Teig
plötzlich übermäßig klebrig ist, dann ist die beste Option, eine Kastenform zu verwenden. Nimm
deinen Teig einfach und gib ihn in eine Kastenform. Warte, bis die Teigmischung
etwas an Volumen zugenommen hat, und backe dann. Du wirst ein sehr leckeres Sauerteigbrot
erhalten. Wenn es etwas zu sauer ist, kannst du deinen Teig einfach länger
backen, um während des Backvorgangs etwas Säure zu verkochen. Du kannst auch
deinen Teig verwenden, um einen neuen Starter anzusetzen und es morgen erneut zu versuchen.
Wenn du hungrig bist, kannst du auch etwas Teig direkt in eine erhitzte Pfanne mit etwas
Öl geben. Das ergibt köstliche Sauerteig-Fladenbrote.

Um Probleme mit Übergärung zu beheben, musst du den Fermentationsprozess
früher stoppen. Was ich~gerne mache, ist eine kleine Fermentationsprobe aus meinem Teig zu entnehmen.
Anhand der Volumenzunahme dieser Probe kann ich~meistens beurteilen, wann meine Fermentation
abgeschlossen ist. Versuche, mit einer Volumenzunahme von \qty{25}{\percent} deines
Hauptteigs oder deiner Probe zu beginnen. Je nachdem, wie viel Gluten dein Mehl hat, kannst du
länger fermentieren. Bei einem starken Mehl mit
\qtyrange{14}{15}{\percent} Protein sollte eine Fermentation
bis zu \qty{100}{\percent} Volumenzunahme problemlos möglich sein. Dies hängt jedoch auch von der
Zusammensetzung deines Sauerteigstarters aus Hefe und Bakterien ab. Je mehr bakterielle Fermentation,
desto schneller baut sich die Teigstruktur ab. Häufiges Füttern deines Sauerteigstarters
verbessert die Hefeaktivität. Außerdem kann ein fester Sauerteigstarter
eine gute Lösung sein. Die verbesserte Hefeaktivität führt zu einem fluffigeren
Teig mit weniger bakterieller Aktivität. Eine bessere Hefeaktivität führt auch
zu weniger Säure im fertigen Brot. Wenn du einen sehr kräftigen, säuerlichen
Geschmack anstrebst, hilft ein stärkeres Mehl mit mehr Gluten.

Beim Retardieren von Sauerteig (Kaltgare im Kühlschrank) spielt die Temperatur eine
entscheidende Rolle bei den Fermentationsraten. Während der Teig im Kühlschrank abkühlt,
verlangsamt sich die Fermentation. Wenn man die Kaltgare bei einer wärmeren
Temperatur beginnt, dauert diese Verlangsamung länger.

Zum Beispiel könnte ein Teig, der nach 8~Stunden Kaltgare ideal ist, bereits in
nur 4~Stunden fertig sein, wenn er bei einer höheren Temperatur gestartet wurde. Daher ist es wichtig,
zu experimentieren und die optimale Kaltgaredauer für deine spezifischen
Bedingungen zu ermitteln. Umgekehrt, wenn der Teig kälter startet, stoppt die Fermentation schneller
im Kühlschrank. In solchen Fällen kann es hilfreich sein, den Teig vor dem Kühlen
kurz bei Raumtemperatur gehen zu lassen.

\subsection{Untergärt}

\begin{figure}
  \centering
  \includegraphics[width=\textwidth]{fermented-too-short-underbaked}
  \caption[Untergärtes Brot]{Ein dichter Teig mit einem gummiartigen, nicht vollständig
      gelatinierten Bereich. Das Bild wurde vom Nutzer
      \emph{wahlfeld} aus unserem Community-Discord-Server bereitgestellt.}%
  \label{fig:fermented-too-short-underbaked}
\end{figure}

Dieser Fehler wird auch häufig als \emph{unterreift} bezeichnet. Allerdings ist unterreift
kein guter Begriff, da er sich nur auf eine zu kurze
Stückgare im Brotbackprozess bezieht.
Wenn man sein Brot nach einer perfekt getimten Stockgare backen würde,
wäre das Ergebnis nicht unterreift, selbst wenn man die Stückgare komplett übersprungen hätte.
Die Stückgare macht den Teig etwas dehnbarer und ermöglicht dem Sauerteig,
den Teig noch etwas mehr aufzublähen. Bei einem untergärten Brot ist
etwas früher schiefgelaufen, während der Stockgare oder vielleicht sogar
vorher beim Sauerteigstarter.

Ein typischer untergärter Teig hat sehr große Lufteinschlüsse und ist stellenweise
feucht und gummiartig. Die großen Einschlüsse lassen sich
mit der Herstellung einer ungesäuerten Weizen- oder Maistortilla vergleichen. Wenn man den Teig in der Pfanne backt,
beginnt das Wasser langsam zu verdampfen. Das Gas wird in der Struktur des Teigs
eingeschlossen und bildet Blasen. Bei einer Tortilla ist das das gewünschte Verhalten.
Beobachtet man diesen Prozess jedoch in einem größeren Teig, entstehen mehrere
Riesenalveolen. Das Wasser verdampft und die ersten Alveolen bilden sich. Dann beginnt irgendwann
die Stärke zu gelatinieren und wird fest. Das geschieht zuerst in den Blasen,
da das Innere schneller erhitzt als der Rest des Teigs. Sobald die gesamte Stärke
gelatiniert ist, behalten die Alveolen ihre Form und dehnen sich nicht weiter aus. Während dieses
Prozesses werden andere Teile des Brotteigs nach außen gedrückt. Deshalb bildet ein untergärter
Teig manchmal sogar ein Ohr beim Backen. Dies
wird auch häufig als \emph{Trugkrume} (engl. fool's crumb) bezeichnet. Man freut sich über ein Ohr, das
ziemlich schwer zu erreichen sein kann. Außerdem denkt man vielleicht, man habe endlich große
Lufteinschlüsse in der Krume erzeugt. Aber in Wirklichkeit hat man zu kurz fermentiert.

\begin{figure}
  \centering
  \includegraphics[width=\textwidth]{fools-crumb}
  \caption[Trugkrume mit großen Alveolen]{Ein typisches Beispiel einer Trugkrume
      mit einem Ohr und mehreren übermäßig großen Alveolen. Das Bild wurde
      von Rochelle aus unserem Community-Discord-Server bereitgestellt.}%
  \label{fig:fools-crumb}
\end{figure}

In einem richtig fermentierten Teig helfen die Alveolen bei der Wärmeübertragung durch den Teig.
Von den vielen kleinen, durch die Fermentation entstandenen Blasen aus gelatiniert die Stärke. Bei
einem untergärten Teig funktioniert diese Wärmeübertragung nicht richtig. Deshalb
gibt es manchmal Bereiche, die wie roher Teig aussehen. Bäcker bezeichnen dies als eine sehr
gummiartige Struktur. Den Teig länger zu backen würde die Stärke in diesen Bereichen ebenfalls
gelatinieren. Allerdings wären dann andere Teile des Brotes zu lange gebacken.

Um Probleme mit Untergärung zu beheben, muss man den Teig einfach
länger fermentieren. Allerdings gibt es eine Obergrenze für die Fermentationszeit,
da das Mehl ab dem Moment des Wasserkontakts abgebaut wird. Deshalb
kann es sinnvoll sein, den Fermentationsprozess einfach zu beschleunigen. Als grobe
Orientierung strebe ich~eine Stockgaredauer von etwa 8--12~Stunden an.
Um das zu erreichen, kann man versuchen, den Sauerteigstarter aktiver zu machen. Das gelingt,
indem man den Starter über mehrere Tage täglich füttert. Verwende dasselbe Verhältnis wie
für den Hauptteig. Wenn du \qty{20}{\percent} Starter
bezogen auf das Mehl verwendest, nimm ein 1:5:5-Verhältnis zum Füttern. Das wären zum Beispiel
\qty{10}{\gram} vorhandener Starter, \qty{50}{\gram} Mehl, \qty{50}{\gram}
Wasser. Um die Hefeaktivität noch weiter zu steigern, kannst du
einen festen Sauerteigstarter
in Betracht ziehen. Die Bakterien produzieren hauptsächlich Säure. Je mehr Säure
sich ansammelt, desto weniger aktiv ist die Hefe. Der feste Sauerteigstarter
ermöglicht es dir, die Fermentation deines Teigs mit stärkerer Hefeaktivität
und weniger bakterieller Aktivität zu beginnen.

\subsection{Zu wenig Teigspannung}

\begin{figure}
  \centering
  \includegraphics[width=\textwidth]{flat-bread}
  \caption{Ein sehr flaches Brot ohne ausreichende Teigspannung.}%
  \label{fig:flat-bread}
\end{figure}

Wenn ein Teig beim Backen stark auseinanderläuft, besteht die Wahrscheinlichkeit,
dass du nicht genug Teigspannung aufgebaut hast. Das bedeutet, dass dein Glutennetzwerk
nicht richtig entwickelt wurde. Dein Teig ist zu dehnbar und flacht eher ab,
anstatt im Ofen nach oben aufzuspringen. Dies kann auch passieren, wenn du
deinen Teig zu lange gehen lässt. Mit der Zeit entspannt sich das Gluten und dein Teig
wird immer dehnbarer. Du kannst dieses Entspannungsverhalten des Glutens
auch beim Pizzabacken beobachten. Direkt nach dem Formen der Teigkugeln ist es
sehr schwer, den Pizzaboden zu formen. Wenn man 30--90~Minuten wartet, wird das Dehnen
des Teigs deutlich einfacher.

Der einfachste Weg, dies zu beheben, ist wahrscheinlich, den Teig am Anfang mehr zu kneten. Zur Vereinfachung
kannst du auch weniger Wasser für dein Mehl verwenden. Das ergibt von Anfang an einen elastischeren Teig.
Dieses Konzept wird häufig bei Sauerteig ohne Kneten angewandt. Alternativ kannst du
auch mehr Dehn- und Faltschritte während der Stockgare durchführen. Jedes
Dehnen und Falten hilft, das Glutennetzwerk zu stärken und einen elastischeren Teig zu erhalten.
Die letzte Möglichkeit, einen Teig mit zu wenig Teigspannung zu korrigieren, ist, den Teig straffer zu formen.

\subsection{Zu heiß gebacken}

\begin{figure}
  \centering
  \includegraphics[width=\textwidth]{baked-too-hot-v2}
  \caption{Ein Brot mit sehr großen Alveolen nahe der Kruste.}%
  \label{fig:baked-too-hot}
\end{figure}

Dies ist ein häufiger Fehler, der mir oft passiert ist. Wenn du deinen Teig
bei zu hoher Temperatur backst, schränkst du die Ausdehnung deines Teigs ein. Die Stärke
gelatiniert und wird immer fester. Bei etwa
\qty{140}{\degreeCelsius} (\qty{284}{\degF}) beginnt die Maillard-Reaktion,
die Kruste deines Brotteigs vollständig zu verdicken. Das ist ähnlich wie wenn man
Brotteig ohne Dampf backt. Während die innere Temperatur des Teigs steigt,
verdampft immer mehr Wasser, Gas dehnt sich aus und der Teig wird nach oben gedrückt.
Sobald der Teig die Kruste erreicht, kann er sich nicht mehr ausdehnen. Die Alveolen verschmelzen
zu größeren Strukturen nahe der Oberfläche des Teigs. Durch zu heißes Backen
erzielst du nicht das Ohr, das zusätzlichen Geschmack verleiht. Außerdem wird durch die eingeschränkte
Ausdehnung die Krume nicht so fluffig, wie sie sein könnte.

Wenn du einen dehnbaren Teig mit hoher Hydratation hast, führt zu kaltes Backen dazu,
dass der Teig stark auseinanderläuft. Die Gelatinierung der Stärke ist
entscheidend dafür, dass der Teig seine Struktur behält. Nach mehreren
Experimenten scheint mein optimaler Wert für maximalen Ofentrieb bei etwa
\qty{230}{\degreeCelsius} (\qty{446}{\degF}) zu liegen. Überprüfe die Temperatur
deines Ofens, denn in einigen
Fällen stimmt die angezeigte Temperatur nicht mit der tatsächlichen Temperatur deines
Ofens überein~\cite{too+hot+baking}. Schalte den Ventilator deines Ofens aus. Die meisten
Haushaltsöfen sind darauf ausgelegt, den Dampf so schnell wie möglich abzuleiten. Wenn du den
Ventilator nicht ausschalten kannst, erwäge die Verwendung eines Bräters (Dutch Oven).

\subsection{Mit zu wenig Dampf gebacken}

\begin{figure}[h]
  \centering
  \includegraphics[width=\textwidth]{no-steam}
  \caption[Brot mit zu wenig Dampf gebacken]{Eines meiner früheren Brote, das
      ich~bei einem Freund gebacken habe, wo ich~den Teig nicht richtig bedampfen konnte.}%
  \label{fig:no-steam}
\end{figure}

Ähnlich wie beim zu heißen Backen bildet sich bei zu wenig Dampf die Kruste deines Teigs
zu schnell. Es ist schwer, den Unterschied zwischen beiden Fehlern zu erkennen.
Ich~frage typischerweise zuerst nach der Temperatur und dann nach der Dampftechnik,
um festzustellen, was beim Backprozess schiefgelaufen sein könnte. Zu wenig Dampf erkennt man
typischerweise an einer dicken Kruste rund um den Teig, gepaart
mit großen Alveolen zu den Rändern hin.

Der Dampf verhindert im Wesentlichen, dass die Maillard-Reaktion zu schnell
an der Kruste einsetzt. Deshalb ist das Bedampfen in den ersten Phasen des Backens so wichtig.
Der Dampf hält die Temperatur der Kruste in der Nähe von etwa
\qty{100}{\degreeCelsius} (\qty{212}{\degF}). Dampf erzeugen
kann man mit einem Bräter (Dutch Oven), einem umgedrehten Blech und/oder einer Schüssel kochendem Wasser.
Möglicherweise hast du auch einen Ofen mit eingebauter Dampffunktion. Alle Methoden funktionieren,
es hängt davon ab, was du zur Hand hast. Meine Standard-Methode ist ein umgedrehtes
Blech auf meinem Teig, zusammen mit einer Schüssel voll kochendem Wasser am Boden
des Ofens.

\begin{figure}[ht]
  \centering
  \includegraphics[width=\textwidth]{apple-experiment-temperatures}
  \caption[Messung von Umgebungs- und Oberflächentemperatur]{Ein Apfel mit 2 Sonden
      zur Messung der Umgebungs- und Oberflächentemperaturen verschiedener Dampftechniken
      in einem Bräter (Dutch Oven).}%
  \label{fig:apple-experiment-temperatures}
\end{figure}

Allerdings kann es auch zu viel Dampf geben. Dafür habe ich~einen Bräter mit großen Eis\-würfeln getestet,
um zusätzlichen Dampf zu erzeugen. Die Temperatur der Teigoberfläche würde direkt
auf fast 100°C springen. Der Dampf enthält mehr Energie und kann daher durch Konvektion
die Oberfläche des Teigs schneller erhitzen. Ich~habe dies getestet, indem ich~einen Apfel in
einen Bräter gelegt und seine Oberflächentemperatur mit einem Grillthermometer gemessen habe.
Dann habe ich~die Dampfmethoden variiert, um darzustellen, wie schnell sich die Temperatur
nahe der Oberfläche ändert. Ich~habe einen Eiswürfel in einem vorgeheizten
Bräter getestet, einen einfachen vorgeheizten Bräter, einen vorgeheizten Bräter mit Wasser\-sprühern
auf der Apfeloberfläche und einen nicht vorgeheizten Bräter, bei dem ich~nur den
Boden vorgeheizt habe. Das Experiment zeigte, dass die Eiswürfel-Methode die
Oberfläche des Apfels deutlich schneller erhitzte. Wenn ich~dies mit einem Brotteig
nachstellte, erzielte ich~weniger Ofentrieb.

\begin{figure}[ht]
  \centering
  \input{plots/fig-temperature-surface.tex}
  \caption[Oberflächentemperatur je nach Dampftechnik]{Eine Grafik, die zeigt, wie
      sich die Temperatur der Apfeloberfläche bei verschiedenen
      Dampftechniken verändert.}%
  \label{fig:apple-experiment-surface-temperatures}
\end{figure}

\begin{figure}[ht]
  \centering
  \begin{tikzpicture}
    \tikzstyle{every node}+=[font=\normalsize\rmfamily]
    \begin{axis}[
        title=Umgebungstemperatur,
        grid=both,
        major grid style={line width=.2pt,draw=gray!30},
        axis x line=middle,
        axis y line=middle,
        axis line style={-Latex},
        width=\textwidth,
        height=0.5\textwidth,
        xmax=35, xmin=-0.1,
        ymax=205, ymin=-0.1,
        every axis y label/.style={%
            at={(ticklabel cs:0.5)},rotate=90,anchor=near ticklabel}, 
        every axis x label/.style={%
            at={(ticklabel cs:0.5)},anchor=near ticklabel}, 
        legend pos = south east,
        legend style={draw=none},
        legend cell align={left},
        xlabel=Dauer (Minuten), ylabel=Temperatur (\SI{}{\degree} C)
        ]
        \addplot [color=redpic,smooth,ultra thick]    table {plots/icecube_ambient.table};
        \addplot [color=codeblack,smooth,ultra thick] table {plots/preheated_ambient.table};
        \addplot [color=codeblue,smooth,ultra thick]  table {plots/non-preheated_ambient.table};
        \addplot [color=yellowpic,smooth,ultra thick] table {plots/preheated_bottom_ambient.table};
        \addplot [color=pinkpic,smooth,ultra thick]   table {plots/spritzing_ambient.table};
        \legend{Eiswürfel, Vorgeheizt, Nicht vorgeheizt, Boden vorgeheizt, Besprühen};
    \end{axis}
\end{tikzpicture}

  \caption[Brätertemperatur je nach Dampftechnik]{Diese Grafik zeigt,
      wie sich die Umgebungstemperaturen im Bräter je nach
      verwendeter Dampftechnik verändern.}%
  \label{fig:apple-experiment-ambient-temperatures}
\end{figure}

Grundsätzlich ist es allerdings relativ schwierig, zu viel Dampf zu erzeugen. Diesen Fehler konnte ich~nur
mit einem Bräter als Dampfmethode in Kombination mit relativ
großen Eiswürfeln produzieren. Nach Gesprächen mit anderen Bäckern, die denselben Bräter verwenden, scheint es,
dass meine Eiswürfel (etwa \qty{80}{\gram}) viermal so schwer waren wie die,
die andere Bäcker verwenden würden (\qty{20}{\gram}).
